\documentclass[11pt,a4paper]{article}
\usepackage[margin=27mm]{geometry}
\usepackage{amsmath,amssymb,amsthm,mathtools}
\usepackage{lmodern,microtype,booktabs}
\usepackage[colorlinks=true,linkcolor=blue,citecolor=blue,urlcolor=blue]{hyperref}
\hypersetup{pdftitle={Mass-only obstructions for Abel-weighted Brownian responses},
pdfauthor={RH--Weil research project}}
\newtheorem{theorem}{Theorem}[section]
\newtheorem{lemma}[theorem]{Lemma}
\newtheorem{corollary}[theorem]{Corollary}
\theoremstyle{definition}
\newtheorem{definition}[theorem]{Definition}
\theoremstyle{remark}
\newtheorem{remark}[theorem]{Remark}
\newcommand{\TV}{\mathrm{TV}}
\newcommand{\ev}{\mathrm{ev}}
\newcommand{\RR}{\mathbb R}
\newcommand{\dd}{\,d}
\DeclareMathOperator{\supp}{supp}
\title{Mass-only obstructions for Abel-weighted Brownian responses}
\author{RH--Weil research project}
\date{September 5, 2026\\[3pt]\small Research draft; literature-priority audit pending}
\begin{document}
\maketitle
\begin{abstract}
We study a two-channel Brownian response built from the actual
von Mangoldt coefficients and a matched continuous Abel source.
For the common cutoff \(N=\lfloor Y(\log Y)^2\rfloor\), an isolated
prime atom gives the unconditional lower bound
\(J_4(Y,N)\gg_\sigma Y^{-9}(\log Y)^{-7}\).
The Mellin transform of the uncut mass discrepancy, together with
a self-contained Landau argument, supplies arbitrarily large mass
zeros whenever the Riemann zeta function has a zero to the right
of \(\Re s=\sigma\). Truncating at those scales leaves a
superpolynomially small mass discrepancy but a polynomially large
response. More generally, a uniform two-parameter lower bound holds
throughout \(Y\le N\le Y^2/8\), extending the obstruction to a class
of regular cutoffs with \(N/(Y\log Y)\to\infty\).
Consequently no positive power of the mass discrepancy
uniformly bounds the response over all sufficiently large real
Abel scales when \(0<\sigma<1/2\). At \(\sigma=1/2\), the proposed
uniform budget would imply RH; no converse is asserted.
The result does not exclude an estimate along a prescribed dyadic
schedule or a selected cofinal subsequence.
\end{abstract}

\section{The response and the main result}

\noindent\textbf{Status convention.}
The theorems and lemmas below are internally proved statements [T];
their no-go consequences are [N]. External classical inputs are [R],
finite computations are [E], and the remaining arithmetic and
literature-priority questions are [O].

Fix \(0<\sigma<1\). Put \(a=1-\sigma\). For \(Y>0\) and
\(1<N\le\infty\), define positive measures on the nonnegative lag axis by
\begin{align}
 \alpha_{Y,N}
 &=\sum_{2\le n\le N}\Lambda(n)n^{-\sigma}e^{-n/Y}\delta_{\log n},
 \label{eq:sources}\\
 d\beta_{Y,N}(\lambda)
 &=\mathbf1_{[0,\log N]}(\lambda)e^{a\lambda-e^\lambda/Y}\dd\lambda.
 \nonumber
\end{align}
The upper endpoint is omitted when \(N=\infty\). Write
\[
 A=\alpha_{Y,N}(\RR),\quad B=\beta_{Y,N}(\RR),\quad
 S=A+B,\quad M=A-B,\quad \mu=M/S.
\]
For a measure on the nonnegative axis let \(\eta^\ev\) denote its even
symmetrization, with half the mass reflected to each side, and set
\[
 p=\alpha_{Y,N}^\ev-A\delta_0,\qquad
 c=B\delta_0-\beta_{Y,N}^\ev,\qquad r=p+c.
\]
For a real signed measure \(\eta\) of total mass zero, let
\(F_\eta(t)=\eta((-\infty,t])\). Our normalization is
\begin{equation}
 \begin{split}
 D&=\|F_p\|_2^2+\|F_c\|_2^2,\\
 P^2&=\|F_{r*r*p}\|_2^2+\|F_{r*r*c}\|_2^2,\qquad
 J_4(Y,N)=\frac{P^2}{S^4D}.
 \end{split}
 \label{eq:response}
\end{equation}
All these quantities are finite for the sources above, and \(S,D>0\).

\begin{theorem}\label{thm:main}
Put \(L=\log Y\) and \(N(Y)=\lfloor YL^2\rfloor\).
\begin{enumerate}
 \item For every fixed \(0<\sigma<1\), uniformly for sufficiently large
 real \(Y\),
 \begin{equation}
 J_4(Y,N(Y))\ge c_\sigma Y^{-9}L^{-7}
 \label{eq:floor}
 \end{equation}
 with \(c_\sigma>0\).
 \item If \(\zeta\) has a nonreal zero \(\rho\) with \(\Re\rho>\sigma\),
 then for every \(\kappa>0\) and every sufficiently large \(Y_0\),
 \begin{equation}
 \sup_{\substack{Y\ge Y_0\\ \mu(Y,N(Y))\ne0}}
 \frac{J_4(Y,N(Y))}{|\mu(Y,N(Y))|^\kappa}=\infty.
 \label{eq:nogo}
 \end{equation}
\end{enumerate}
\end{theorem}

\begin{corollary}\label{cor:rh}
Conclusion \eqref{eq:nogo} holds unconditionally for \(0<\sigma<1/2\).
For any \(0<\sigma<1\), an eventual bound
\(J_4(Y,N(Y))\le C|\mu(Y,N(Y))|^4\) on all real scales with nonzero
mass discrepancy implies that \(\zeta\) has no zero in \(\Re s>\sigma\).
In particular, at \(\sigma=1/2\) it implies RH.
\end{corollary}

Only the existence of a critical-line zero is needed for the first
assertion of the corollary; this is a classical theorem
\cite{DLMF}. The no-go concerns the mass-only budget in
\eqref{eq:response}. It is neither a disproof of RH nor an obstruction
to every realization of a Weil-type configuration.

\section{Brownian energy and a positive actual response}

We use the Fourier convention
\(\widehat\eta(\xi)=\int_{\RR}e^{-i\xi x}\dd\eta(x)\).

\begin{lemma}\label{lem:energy}
For a real signed measure \(\eta\) of mass zero and finite first
absolute moment,
\begin{equation}
 \|F_\eta\|_2^2
 =-\frac12\iint |x-y|\dd\eta(x)\dd\eta(y)
 =\frac1{2\pi}\int_\RR\frac{|\widehat\eta(\xi)|^2}{\xi^2}\dd\xi.
 \label{eq:energy}
\end{equation}
For \(k_\lambda=(\delta_\lambda+\delta_{-\lambda})/2-\delta_0\),
\[
 \langle F_{k_\lambda},F_{k_\gamma}\rangle
 =\tfrac12\min(\lambda,\gamma)\qquad(\lambda,\gamma\ge0).
\]
\end{lemma}
\begin{proof}
The squared difference of two half-line indicators has integral
\(|x-y|\). Expanding its integral against \(\eta\otimes\eta\), the two
single-variable terms vanish because \(\eta(\RR)=0\).
The Fourier formula follows from \(F_\eta'=\eta\) and Plancherel.
Finite first absolute moment justifies the distance integral.
For the final identity, each primitive is constant with magnitude
\(1/2\) on each half of \([-\lambda,\lambda]\), with opposite signs
on the two halves. Their overlap has length \(2\min(\lambda,\gamma)\).
\end{proof}

\begin{lemma}\label{lem:positive}
For every \(Y>0\), the uncut response \(J_4(Y,\infty)\) is strictly
positive and continuous in \(Y\).
\end{lemma}
\begin{proof}
The measure \(r\) has a nonzero atom at \(\log2\). The continuous
source has no atom there, so \(r\ne0\).
Fourier uniqueness gives some \(\xi\ne0\) with \(\widehat r(\xi)\ne0\).
For every real \(\xi\ne0\),
\[
 \widehat c(\xi)=\int_1^\infty x^{-\sigma}e^{-x/Y}
                    (1-\cos(\xi\log x))\dd x>0.
\]
Thus \(\widehat{r*r*c}=\widehat r^{\,2}\widehat c\) is not identically
zero, and Lemma~\ref{lem:energy} gives strictly positive energy.
On any compact \(Y\)-interval the source coefficients and their
first absolute logarithmic moments admit a common integrable
majorant. They vary continuously in the corresponding weighted
total-variation norm. Convolution and the distance integral preserve
this continuity, proving the assertion.
\end{proof}

\section{Mellin poles force mass cancellation}

Write
\[
 M_\infty(Y)=
 \sum_{n\ge2}\Lambda(n)n^{-\sigma}e^{-n/Y}
 -\int_1^\infty x^{-\sigma}e^{-x/Y}\dd x.
\]
\begin{lemma}\label{lem:mellin}
For \(\Re s>1-\sigma\),
\begin{equation}
 \int_0^\infty M_\infty(Y)Y^{-s-1}\dd Y
 =\Gamma(s)\left(-\frac{\zeta'}{\zeta}(\sigma+s)
                       -\frac1{\sigma+s-1}\right).
 \label{eq:mellin}
\end{equation}
The meromorphic function on the right is holomorphic near every
positive real point. At \(s=\rho-\sigma\), for a nonreal zero \(\rho\)
of multiplicity \(m_\rho\), its residue is
\(-m_\rho\Gamma(\rho-\sigma)\ne0\) whenever \(\Re\rho>\sigma\).
\end{lemma}
\begin{proof}
The bound \(\Lambda(n)\le\log n\) gives
\(A_\infty+B_\infty\ll_\sigma Y^{1-\sigma}(1+\log Y)\) for \(Y\ge2\);
at zero both masses decay exponentially in \(1/Y\).
Absolute convergence therefore permits interchange of integrals and
sums. The substitution \(t=x/Y\) gives
\(\int_0^\infty e^{-x/Y}Y^{-s-1}\dd Y=\Gamma(s)x^{-s}\).
The Euler logarithmic derivative and the continuous integral yield
\eqref{eq:mellin}.

There are no positive real zeros of \(\zeta\): for \(w>1\) use its
Euler product, and for \(0<w<1\) use
\[
 \zeta(w)=\frac{\sum_{n\ge1}(-1)^{n-1}n^{-w}}{1-2^{1-w}}<0.
\]
At \(w=1\) the two displayed pole parts cancel. The Gamma function
is holomorphic on the positive real axis. At a nonreal zero,
the negative logarithmic derivative has residue \(-m_\rho\), and Gamma has
no zero, establishing the last assertion.
\end{proof}

We include the classical positivity argument to specify the analytic
continuation required in the proof. A modern statement of this
Landau mechanism appears in \cite[Theorem 3.1]{MM}; we make no
novelty claim for the mechanism.
\begin{lemma}[Landau mechanism]\label{lem:landau}
Suppose \(f(Y)\ge0\) for \(Y\ge Y_0\ge1\), and the Mellin integral
\(G(s)=\int_{Y_0}^\infty f(Y)Y^{-s-1}\dd Y\) has finite convergence
abscissa \(c\). Then \(G\) cannot continue holomorphically across
the real point \(c\).
\end{lemma}
\begin{proof}
If it could, choose a small \(\varepsilon>0\) such that the Taylor
radius at \(b=c+\varepsilon\) exceeds \(\varepsilon\).
For every integer \(k\ge0\),
\[
 (-1)^kG^{(k)}(b)
 =\int_{Y_0}^\infty f(Y)(\log Y)^kY^{-b-1}\dd Y\ge0.
\]
Choose \(h>\varepsilon\) within that Taylor radius. Taylor convergence
and Tonelli imply
\[
 \int_{Y_0}^\infty f(Y)Y^{-(b-h)-1}\dd Y
 =\sum_{k\ge0}\frac{h^k}{k!}(-1)^kG^{(k)}(b)<\infty,
\]
contrary to \(b-h<c\).
\end{proof}

\begin{theorem}\label{thm:oscillation}
If \(\zeta\) has a nonreal zero \(\rho\) with \(\Re\rho>\sigma\),
then \(M_\infty\) takes both strictly positive and strictly negative
values arbitrarily far out. In particular, it has a sequence of
zeros \(Y_j\to\infty\).
\end{theorem}
\begin{proof}
Assume first that \(M_\infty\) is eventually nonnegative.
Its nonnegative tail Mellin transform has convergence abscissa
\(c\le1-\sigma\). Removing a bounded initial \(Y\)-interval from
\eqref{eq:mellin} subtracts an entire function, because of the
exponential decay at zero. If \(c<\Re\rho-\sigma\), the tail
integral would be holomorphic at \(\rho-\sigma\), contradicting
the pole there and uniqueness of meromorphic continuation.
Hence \(c\ge\Re\rho-\sigma>0\).
Lemma~\ref{lem:landau} now contradicts holomorphy at the positive
real point \(c\). Apply the same argument to \(-M_\infty\) to
exclude eventual nonpositivity.
The mass is real analytic for \(Y>0\), by locally dominated
differentiation of its defining sums and integrals. The intermediate
value theorem supplies the stated zeros.
\end{proof}

By Lemma~\ref{lem:positive}, even omitting all exact mass zeros cannot
rescue a uniform uncut mass-only bound: the response stays positive
while the mass tends to zero near each root.
The next section proves a quantitative version that survives a
finite arithmetic cutoff.

\section{An isolated prime atom survives the continuous background}

\begin{proof}[Proof of Theorem~\ref{thm:main}(1)]
Take \(L=\log Y\), \(N=\lfloor YL^2\rfloor\).
Bertrand's theorem \cite{Erdos} provides a prime \(q\in(Y,2Y)\);
apply its integer form at \(\lfloor Y\rfloor\).
For sufficiently large \(Y\), we have \(q\le N<q^2\).
Put \(d=\alpha^\ev-M\delta_0\) and \(b=\beta^\ev\), so \(r=d-b\).
The atomic part of \(\eta=r*r*p\) is \(d*d*p\). Its remaining part
\[
 g(x)\dd x=-2d*b*p+b*b*p
\]
is absolutely continuous.

At \(x_0=3\log q\) the atomic coefficient is exactly
\begin{equation}
 w=\frac18\bigl((\log q)q^{-\sigma}e^{-q/Y}\bigr)^3
 \gg_\sigma L^3Y^{-3\sigma}.
 \label{eq:jump}
\end{equation}
Indeed every atomic term uses at most three signed prime-power
lags. Since \(q^2>N\), each lag contributes at most one to the
\(q\)-adic valuation. A total valuation of three requires three
positive lags divisible by \(q\); their product is \(q^3\), forcing
each factor to equal \(q\). Neither a center term nor a negative
lag can participate. An absolutely continuous term cannot alter
this atomic coefficient.

Every other atomic location is \(\log(u/v)\) for positive integers
\(u,v\le N^3\). If \(u/v\ne q^3\), the distinct integers \(u,vq^3\)
have smaller member at most \(N^3\), and therefore
\begin{equation}
 |\log(u/v)-3\log q|
 \ge\log(1+N^{-3})\ge(2N^3)^{-1}.
 \label{eq:spacing}
\end{equation}
Set \(\delta=(8N^3)^{-1}\), which isolates the atom on both sides.

Elementary summation and maximization of the continuous density give
\[
 S\ll_\sigma Y^{1-\sigma}L,\qquad
 \left\|\frac{db}{dx}\right\|_\infty\ll_\sigma Y^{1-\sigma}.
\]
Here \(x\) denotes the additive lag coordinate. Since
\(\|d\|_\TV\le2S\), \(\|p\|_\TV\le2A\), and \(\|b\|_\TV=B\),
\begin{equation}
 \|g\|_\infty
 \le10S^2\left\|\frac{db}{dx}\right\|_\infty
 \ll_\sigma Y^{3(1-\sigma)}L^2.
 \label{eq:density}
\end{equation}
Equations \eqref{eq:jump}--\eqref{eq:density} imply
\(\delta\|g\|_\infty/w\ll_\sigma L^{-7}\), eventually at most \(1/4\).

The two one-sided limits of \(F_\eta\) at \(x_0\) differ by \(w\);
one therefore has absolute value at least \(w/2\).
On the corresponding adjacent interval of length \(\delta\) there
are no other atoms, and the continuous change is at most \(w/4\).
Consequently
\[
 P^2\ge\|F_\eta\|_2^2\ge\delta w^2/16
 \gg_\sigma Y^{-6\sigma-3}.
\]
Lemma~\ref{lem:energy} also gives
\[
 D\le\tfrac12(\log N)(A^2+B^2),\qquad
 S^4D\ll_\sigma Y^{6(1-\sigma)}L^7.
\]
Dividing yields \eqref{eq:floor}. This proof permits the jumps of
the integer-valued cutoff \(N(Y)\).
\end{proof}

\section{Finite-cutoff mass cancellation and the obstruction}

\begin{proof}[Proof of Theorem~\ref{thm:main}(2)]
Choose the uncut mass zeros \(Y_j\) from
Theorem~\ref{thm:oscillation}, and set \(L_j=\log Y_j\),
\(N_j=\lfloor Y_jL_j^2\rfloor\).
At each such point the finite mass is the negative of a difference
of omitted tails, so
\begin{equation}
 |M(Y_j,N_j)|\ll_\sigma
 Y_j^{1-\sigma}L_j^2e^{-L_j^2/2}.
 \label{eq:tail}
\end{equation}
For completeness, bound the prime tail by the integral of
\((\log x)x^{-\sigma}e^{-x/Y}\) from \(N\), which is decreasing
there. After \(x=Yu\), use \(u^{-\sigma}\le1\) and \(\log u\le u\)
for \(u\ge1\). The resulting upper bound is
\[
 C_\sigma Y^{1-\sigma}(L+N/Y+1)e^{-N/Y}.
\]
The continuous tail is smaller. This proves \eqref{eq:tail};
the floor error in \(N/Y\) is harmless.
The finite continuous mass on \([Y/2,Y]\) is
\(\gg_\sigma Y^{1-\sigma}\). Hence
\[
 |\mu(Y_j,N_j)|\ll_\sigma L_j^2e^{-L_j^2/2}.
\]

To ensure nonzero finite mass, perturb \(Y_j\) slightly to the right.
The cutoff is constant on a right neighborhood, even when \(Y_j\)
is a cutoff threshold. For fixed \(N_j\), the function
\(M(Y,N_j)\) is real analytic and not identically zero: at small
\(Y\) its integral beginning at \(1\) dominates the prime terms
beginning at \(2\). Its zeros are therefore isolated.
Choose \(0<Y'_j-Y_j<Y_j^{-1}\) within the same cutoff cell, with
\(M(Y'_j,N_j)\ne0\), and small enough to retain
\eqref{eq:tail} up to a constant.

Theorem~\ref{thm:main}(1) applies directly to these finite sources;
no response-tail transfer is required. For \(\kappa=4\),
\[
 \frac{J_4(Y'_j,N_j)}{|\mu(Y'_j,N_j)|^4}
 \gg_\sigma e^{2L_j^2-9L_j}L_j^{-15}\longrightarrow\infty.
\]
For general \(\kappa>0\), the exponential factor is
\(e^{\kappa L_j^2/2-9L_j}\), which still dominates all logarithmic
factors. This proves \eqref{eq:nogo}.
\end{proof}

\begin{proof}[Proof of Corollary~\ref{cor:rh}]
For \(0<\sigma<1/2\), a classical nonreal critical-line zero has
real part exceeding \(\sigma\), so apply Theorem~\ref{thm:main}(2).
In general, any zero to the right of \(\sigma\) contradicts the
proposed uniform budget. At \(\sigma=1/2\), the symmetry of the
nontrivial zeros about the critical line then gives RH.
\end{proof}

\section{Uniformity across a class of cutoff rules}

\begin{theorem}\label{thm:uniformcutoff}
For every fixed \(0<\sigma<1\), there are \(c_\sigma>0,Y_\sigma\)
such that, simultaneously for every real \(Y\ge Y_\sigma\) and
every integer \(Y\le N\le Y^2/8\),
\begin{equation}
 J_4(Y,N)\ge c_\sigma Y^{-6}N^{-3}(\log Y)^{-1}.
 \label{eq:uniformfloor}
\end{equation}
The constant and the threshold are independent of \(N\).
\end{theorem}
\begin{proof}
This time apply Bertrand's theorem at \(\lfloor Y/2\rfloor\) to
choose \(Y/2<q<Y\). Then \(q\le N<q^2\). The valuation argument
and integer separation from Section 4 apply without change:
the atom at \(3\log q\) has coefficient
\(w\gg_\sigma L^3Y^{-3\sigma}\) and can be isolated with
\(\delta=(8N^3)^{-1}\).
The uncut source bounds give, uniformly in the allowed \(N\),
\[
 S\ll_\sigma Y^{1-\sigma}L,\qquad
 \|g\|_\infty\ll_\sigma Y^{3(1-\sigma)}L^2,\qquad
 \frac{\delta\|g\|_\infty}{w}
 \ll_\sigma\frac{Y^3}{N^3L}\le \frac{C_\sigma}{L}.
\]
For a threshold independent of \(N\), the jump argument therefore
gives
\[
 P^2\ge\delta w^2/16
 \gg_\sigma N^{-3}L^6Y^{-6\sigma}.
\]
Since \(\log N\le2L\), the denominator satisfies
\(S^4D\ll_\sigma Y^{6(1-\sigma)}L^7\).
Dividing proves \eqref{eq:uniformfloor}.
Only \(\log N\asymp\log Y\), not asymptotic equality, is used.
\end{proof}

\begin{theorem}\label{thm:cutoffclass}
Let \(\Phi\) be continuous and nondecreasing on a final real
interval, and put \(N(Y)=\lfloor\Phi(Y)\rfloor\).
Assume eventually that
\[
 Y\le N(Y)\le Y^2/8,\qquad
 \frac{N(Y)}{Y\log Y}\longrightarrow\infty.
\]
If zeta has a nonreal zero to the right of \(\sigma\), then
\eqref{eq:nogo} holds for this cutoff rule, for every fixed
\(\kappa>0\) and every sufficiently large \(Y_0\).
\end{theorem}
\begin{proof}
Take the uncut mass roots \(Y_j\) from
Theorem~\ref{thm:oscillation}, and put \(h_j=N(Y_j)/Y_j\).
The same decreasing-integral estimate used in
\eqref{eq:tail}, now with lower limit \(h_j\), gives
\[
 |M(Y_j,N_j)|\ll_\sigma
 Y_j^{1-\sigma}(L_j+h_j+1)e^{-h_j}.
\]
The continuum mass on \([Y_j/2,Y_j]\) bounds \(S\) below.
Since \(h_j/L_j\to\infty\), this yields
\[
 |\mu(Y_j,N_j)|\ll_\sigma
 (L_j+h_j+1)e^{-h_j}\le e^{-h_j/2}
\]
eventually.
Continuity and monotonicity of \(\Phi\) imply that
\(N(Y)=N_j\) in a right neighborhood of \(Y_j\): continuity
keeps \(\Phi(Y)<N_j+1\), and monotonicity keeps
\(\Phi(Y)\ge N_j\).
This includes the case of an integer value \(\Phi(Y_j)\).
The fixed-\(N_j\) analytic perturbation from Section 5 therefore
provides nonzero mass points \(Y'_j\), with
\(0<Y'_j-Y_j<Y_j^{-1}\), retaining the bound up to a factor of two.

Theorem~\ref{thm:uniformcutoff} and \(N_j\le(Y'_j)^2/8\) give
\(J_4(Y'_j,N_j)\gg_\sigma(Y'_j)^{-12}(\log Y'_j)^{-1}\).
Consequently
\[
 \frac{J_4(Y'_j,N_j)}{|\mu(Y'_j,N_j)|^\kappa}
 \gg_{\sigma,\kappa}
 L_j^{-1}\exp(\kappa h_j/2-12L_j)\longrightarrow\infty.
\]
\end{proof}

This includes each fixed cutoff
\(N=\lfloor Y(\log Y)^b\rfloor\) with \(b>1\), and
\(N=\lfloor Y^{1+\varepsilon}\rfloor\) with \(0<\varepsilon<1\).
For any selected discrete schedule within \(Y\le N\le Y^2/8\),
the budget \(J_4\le C|\mu|^\kappa\) at nonzero mass points
necessarily requires
\[
 |\mu|\ge(c_\sigma/C)^{1/\kappa}
 Y^{-6/\kappa}N^{-3/\kappa}(\log Y)^{-1/\kappa}.
\]
This is not sufficient for the budget, and no violation of this
necessary condition on a prescribed dyadic schedule is proved.

\section{Scope, the Weil interface, and reproducibility}

The continuous-scale quantifier in the no-go theorems is essential.
A bound for \(Y=\theta2^m\), simultaneously with the same constant
and starting index for all \(1\le\theta<2\), would cover all large
real scales and is subject to the theorem.
A prescribed \(\theta\), constants depending on \(\theta\), or an
existentially chosen cofinal schedule are not covered.
In particular, a mass-relative estimate on a selected dyadic
schedule remains an open arithmetic question.

The response \eqref{eq:response} is the two-channel prime/continuous
subsystem used in a conditional capture and Schur-gain criterion.
It is not the complete Weil quadratic form. Incorporating the Gamma
response, or constructing a cohomological realization, requires a
separate bridge. The theorem only prevents treating a uniform
continuous-scale mass-only budget as automatic positivity.
The implication to RH at the central parameter is necessary, not
an equivalence established here.

The argument uses the classical Euler identity, analytic continuation
of zeta, existence of a critical-line zero, and Bertrand's theorem.
It uses no PNT error estimate or assumption about unknown zeros.
For other \(L\)-functions, positive real zeros or complex coefficients
can invalidate the real-axis or positivity steps; no unrestricted
GRH extension is claimed.

The accompanying exact script
\path{scripts/abel_prime_atom_audit.py} checks multiplicative
convolution coefficients and integer separation using rational
arithmetic. The numerical script
\path{scripts/abel_mass_discrepancy_probe.py} checks bounded actual
Abel samples using numpy and mpmath. All samples at
\(\sigma=1/4,1/2,3/4\), \(Y=2^3,\ldots,2^{18}\) were negative.
This neither contradicts eventual oscillation nor estimates its first
occurrence. The computations are not inputs to the asymptotic proofs.
On the Windows test platform, numpy longdouble has the same
52 explicit mantissa bits as binary64. The printed tail majorants
do not include floating-point coefficient or summation errors and
are not certified total-error bounds.
Independent internal audits checked the Mellin singularities,
prime-atom lower bound, cutoff transfer, and scale quantifiers.
The literature priority of their combination remains to be assessed.

\begin{thebibliography}{9}
\bibitem{DLMF}
F.~W.~J. Olver et al., eds.,
\emph{NIST Digital Library of Mathematical Functions},
Section 25.10, Zeros of the Riemann zeta function.
\url{https://dlmf.nist.gov/25.10}.

\bibitem{Erdos}
P.~Erd\H{o}s,
Beweis eines Satzes von Tschebyschef,
\emph{Acta Litt. Sci. Szeged} \textbf{5} (1932), 194--198.
\url{https://users.renyi.hu/~p_erdos/1932-01.pdf}.

\bibitem{MM}
K.~Mahatab and A.~Mukhopadhyay,
\emph{Measure Theoretic Aspects of Oscillations of Error Terms},
arXiv:1512.03144v4 (2018), Theorem 3.1.
\url{https://arxiv.org/abs/1512.03144}.
\end{thebibliography}
\end{document}
